% !TeX spellcheck = ru_RU
% !TEX root = vkr.tex

\section{Математические основы} \label{sec:background}

В настоящем разделе вводятся формальные определения и математический аппарат,
необходимые для строгой постановки задачи классификации инструментов, оценки
качества моделей и статистического сравнения результатов.

%% ===================================================================
\subsection{Формализация задачи бинарной классификации}
\label{subsec:binary-classification}

Пусть $\mathcal{X}$ --- пространство признаков, элементами которого являются
текстовые запросы пользователей к AI-агенту. Определим множество меток
$\mathcal{Y} = \{0, 1\}$, где метка $0$ соответствует классу \texttt{TAG}
(тегирование), а метка $1$ --- классу \texttt{GEO} (геокодирование).

\begin{definition}[Классификатор]
\label{def:classifier}
Классификатором называется измеримое отображение
$f_\theta \colon \mathcal{X} \to \mathcal{Y}$,
параметризованное вектором параметров $\theta \in \Theta$, где
$\Theta$ --- пространство параметров модели.
\end{definition}

Пусть дана выборка $\{(x_i, y_i)\}_{i=1}^{n}$, где
$x_i \in \mathcal{X}$, $y_i \in \mathcal{Y}$.
Обозначим предсказания классификатора $\hat{y}_i = f_\theta(x_i)$.

\begin{definition}[Матрица ошибок]
\label{def:confusion-matrix}
Матрицей ошибок (confusion matrix) классификатора $f_\theta$ на выборке $\{(x_i, y_i)\}_{i=1}^{n}$ называется набор из четырёх величин:
\begin{align}
    \mathrm{TP} &= \bigl|\{i : \hat{y}_i = 1 \;\wedge\; y_i = 1\}\bigr|, \label{eq:tp} \\
    \mathrm{FP} &= \bigl|\{i : \hat{y}_i = 1 \;\wedge\; y_i = 0\}\bigr|, \label{eq:fp} \\
    \mathrm{TN} &= \bigl|\{i : \hat{y}_i = 0 \;\wedge\; y_i = 0\}\bigr|, \label{eq:tn} \\
    \mathrm{FN} &= \bigl|\{i : \hat{y}_i = 0 \;\wedge\; y_i = 1\}\bigr|, \label{eq:fn}
\end{align}
где $|\cdot|$ обозначает мощность множества. Очевидно,
$\mathrm{TP} + \mathrm{FP} + \mathrm{TN} + \mathrm{FN} = n$.
\end{definition}

Положительным классом в дальнейшем считается \texttt{GEO} ($y = 1$);
предметная интерпретация двух типов ошибок (TAG\,$\to$\,GEO и GEO\,$\to$\,TAG)
для рассматриваемого AI-агента приведена в разделе~\ref{sec:task}.


%% ===================================================================
\subsection{Метрики качества классификации}
\label{subsec:metrics}

Пусть $\mathcal{C} = \{0, 1\}$ --- множество классов,
$|\mathcal{C}| = 2$. Для каждого класса $c \in \mathcal{C}$ обозначим
через $\mathrm{TP}_c$, $\mathrm{FP}_c$, $\mathrm{TN}_c$, $\mathrm{FN}_c$
элементы матрицы ошибок, построенной в стратегии <<один против всех>>
(one-vs-rest) для класса $c$.

\begin{definition}[Accuracy и error rate]
\label{def:accuracy}
Доля правильных ответов (accuracy) классификатора определяется как
\begin{equation}
    \mathrm{Acc} = \frac{\mathrm{TP} + \mathrm{TN}}
                        {\mathrm{TP} + \mathrm{TN} + \mathrm{FP} + \mathrm{FN}}.
    \label{eq:accuracy}
\end{equation}
Дополнительной величиной, эквивалентной accuracy, но удобной для оценки
прогресса в режиме <<малого числа ошибок>>, является \emph{доля ошибок}
(error rate):
\begin{equation}
    \mathrm{ER} \;=\; 1 - \mathrm{Acc}
                  \;=\; \frac{\mathrm{FP} + \mathrm{FN}}
                             {\mathrm{TP} + \mathrm{TN} + \mathrm{FP} + \mathrm{FN}}.
    \label{eq:error-rate}
\end{equation}
\end{definition}

\begin{definition}[Precision]
\label{def:precision}
Точность (precision) для класса $c \in \mathcal{C}$ определяется как
\begin{equation}
    P_c = \frac{\mathrm{TP}_c}{\mathrm{TP}_c + \mathrm{FP}_c}.
    \label{eq:precision}
\end{equation}
Величина $P_c$ характеризует долю истинных представителей класса $c$
среди всех объектов, отнесённых классификатором к классу $c$.
\end{definition}

\begin{definition}[Recall]
\label{def:recall}
Полнота (recall) для класса $c \in \mathcal{C}$ определяется как
\begin{equation}
    R_c = \frac{\mathrm{TP}_c}{\mathrm{TP}_c + \mathrm{FN}_c}.
    \label{eq:recall}
\end{equation}
Величина $R_c$ характеризует долю верно распознанных представителей
класса $c$ среди всех объектов, истинно принадлежащих классу $c$.
\end{definition}

\begin{definition}[$F_1$-мера]
\label{def:f1}
$F_1$-мера для класса $c$ определяется как среднее гармоническое точности
и полноты:
\begin{equation}
    F_{1,c} = \frac{2\,P_c\,R_c}{P_c + R_c}
            = \frac{2\,\mathrm{TP}_c}{2\,\mathrm{TP}_c + \mathrm{FP}_c + \mathrm{FN}_c}.
    \label{eq:f1}
\end{equation}
\end{definition}

Для обобщения метрик на многоклассовый случай (а также для единообразного
представления результатов бинарной классификации) используются агрегированные
варианты $F_1$-меры.

\begin{definition}[Macro $F_1$]
\label{def:macro-f1}
Макро-усреднённая $F_1$-мера определяется как среднее арифметическое
$F_1$-мер по всем классам:
\begin{equation}
    F_1^{\mathrm{macro}} = \frac{1}{|\mathcal{C}|}
        \sum_{c \in \mathcal{C}} F_{1,c}.
    \label{eq:macro-f1}
\end{equation}
Данная метрика придаёт равный вес каждому классу вне зависимости от его
представленности в выборке.
\end{definition}

\begin{definition}[Weighted $F_1$]
\label{def:weighted-f1}
Взвешенная $F_1$-мера определяется как
\begin{equation}
    F_1^{\mathrm{weighted}} = \sum_{c \in \mathcal{C}}
        \frac{n_c}{N}\,F_{1,c},
    \label{eq:weighted-f1}
\end{equation}
где $n_c = |\{i : y_i = c\}|$ --- число объектов класса $c$ в выборке,
$N = \sum_{c \in \mathcal{C}} n_c$ --- общий объём выборки.
Эта метрика учитывает дисбаланс классов.
\end{definition}

Систематический анализ метрик классификации и их свойств приведён
в~\cite{sokolova2009systematic}.


%% ===================================================================
\subsection{Низкоранговая адаптация (LoRA)}
\label{subsec:lora}

Метод низкоранговой адаптации (Low-Rank Adaptation, LoRA)~\cite{hu2022lora}
позволяет эффективно дообучать предобученные языковые модели, модифицируя лишь
малую часть параметров.

Пусть $W_0 \in \mathbb{R}^{d \times k}$ --- матрица весов некоторого
линейного слоя предобученной модели. В процессе полного дообучения
(full fine-tuning) вычисляется обновление $\Delta W \in \mathbb{R}^{d \times k}$,
после чего используется матрица $W = W_0 + \Delta W$.
Количество обучаемых параметров при этом составляет $d \cdot k$.

Ключевая идея LoRA состоит в параметризации обновления через произведение
двух матриц малого ранга:
\begin{equation}
    W = W_0 + \Delta W = W_0 + B\,A,
    \label{eq:lora-decomposition}
\end{equation}
где $B \in \mathbb{R}^{d \times r}$, $A \in \mathbb{R}^{r \times k}$,
а $r \ll \min(d, k)$ --- ранг адаптации.

\begin{definition}[Низкоранговая адаптация]
\label{def:lora}
Низкоранговой адаптацией с рангом $r$ называется замена обновления
$\Delta W \in \mathbb{R}^{d \times k}$ на произведение
$\Delta W = B\,A$, где $B \in \mathbb{R}^{d \times r}$,
$A \in \mathbb{R}^{r \times k}$, при фиксированной исходной матрице $W_0$.
\end{definition}

Количество обучаемых параметров при использовании LoRA для одного слоя
составляет
\begin{equation}
    r\,(d + k)
    \label{eq:lora-params}
\end{equation}
вместо $d \cdot k$. При $r \ll \min(d, k)$ достигается значительное
сокращение числа параметров:
\begin{equation}
    \frac{r\,(d + k)}{d \cdot k}
    = r\left(\frac{1}{k} + \frac{1}{d}\right)
    \ll 1.
    \label{eq:lora-ratio}
\end{equation}

Инициализация матриц выполняется следующим образом:
\begin{equation}
    A_{ij} \sim \mathcal{N}(0, \sigma^2), \qquad B = \mathbf{0}_{d \times r},
    \label{eq:lora-init}
\end{equation}
что гарантирует $\Delta W = BA = \mathbf{0}$ на момент начала обучения.
Таким образом, дообучение начинается с исходной предобученной модели без
внесения начальных возмущений.

Прямой проход через адаптированный слой для входного вектора
$h \in \mathbb{R}^k$ вычисляется как
\begin{equation}
    h' = W_0\,h + \frac{\alpha}{r}\,B\,A\,h,
    \label{eq:lora-forward}
\end{equation}
где $\alpha > 0$ --- масштабирующий гиперпараметр, контролирующий вклад
адаптации. Множитель $\alpha / r$ обеспечивает стабильность обучения
при изменении ранга $r$.


%% ===================================================================
\subsection{Статистические критерии сравнения классификаторов}
\label{subsec:statistical-tests}

Для статистически обоснованного сравнения нескольких классификаторов,
обученных на одних и тех же данных и оценённых на общей тестовой выборке,
необходимо применять специализированные критерии, учитывающие зависимость
предсказаний.

%% -------------------------------------------------------------------
\subsubsection{Критерий $Q$ Кохрена}
\label{subsubsec:cochran}

Критерий $Q$ Кохрена~\cite{cochran1950comparison} применяется для проверки
гипотезы об однородности $k \geqslant 2$ бинарных классификаторов,
оценённых на одной и той же выборке из $n$ объектов.

Пусть $X_{ij} \in \{0, 1\}$ --- индикатор правильного ответа $j$-го
классификатора ($j = 1, \ldots, k$) на $i$-м объекте ($i = 1, \ldots, n$).
Введём обозначения:
\begin{equation}
    T_j = \sum_{i=1}^{n} X_{ij}, \qquad
    R_i = \sum_{j=1}^{k} X_{ij}, \qquad
    \bar{T} = \frac{1}{k}\sum_{j=1}^{k} T_j.
    \label{eq:cochran-notation}
\end{equation}
Здесь $T_j$ --- количество правильных ответов $j$-го классификатора,
$R_i$ --- количество классификаторов, давших правильный ответ на $i$-й объект.

\begin{definition}[Нулевая гипотеза критерия Кохрена]
\label{def:cochran-h0}
$H_0$: все $k$ классификаторов имеют одинаковую вероятность правильного
ответа, т.\,е.
$P(X_{i1} = 1) = P(X_{i2} = 1) = \ldots = P(X_{ik} = 1)$
для всех $i$.
\end{definition}

Статистика критерия Кохрена определяется формулой
\begin{equation}
    Q = \frac{k(k-1)\displaystyle\sum_{j=1}^{k}(T_j - \bar{T})^{2}}
             {k\displaystyle\sum_{i=1}^{n}R_i
              - \displaystyle\sum_{i=1}^{n}R_i^{2}}.
    \label{eq:cochran-Q}
\end{equation}

При справедливости $H_0$ и достаточно большом $n$ статистика $Q$ имеет
асимптотическое распределение хи-квадрат:
\begin{equation}
    Q \xrightarrow{d} \chi^{2}(k - 1).
    \label{eq:cochran-distribution}
\end{equation}

Гипотеза $H_0$ отвергается при $Q > \chi^{2}_{1-\alpha}(k-1)$, где
$\chi^{2}_{1-\alpha}(k-1)$ --- $(1-\alpha)$-квантиль распределения
$\chi^{2}$ с $k-1$ степенями свободы. Отвержение $H_0$ означает, что
хотя бы один классификатор статистически значимо отличается от остальных
по доле правильных ответов.

%% -------------------------------------------------------------------
\subsubsection{Критерий Макнемара}
\label{subsubsec:mcnemar}

Для попарного сравнения двух классификаторов $A$ и $B$, оценённых на
одной и той же выборке из $n$ объектов, применяется критерий
Макнемара~\cite{mcnemar1947note, dietterich1998approximate}.

Построим таблицу сопряжённости $2 \times 2$ по результатам работы
двух классификаторов:
\begin{equation}
    \begin{array}{c|cc}
          & B \text{ верно} & B \text{ неверно} \\ \hline
    A \text{ верно}   & n_{11} & n_{10} \\
    A \text{ неверно} & n_{01} & n_{00}
    \end{array}
    \label{eq:mcnemar-table}
\end{equation}
где $n_{10}$ --- количество объектов, на которых $A$ ответил верно,
а $B$ --- неверно; $n_{01}$ --- количество объектов, на которых $B$
ответил верно, а $A$ --- неверно. Величины $n_{11}$ и $n_{00}$
(согласованные ответы) не несут информации о различии классификаторов.

\begin{definition}[Нулевая гипотеза критерия Макнемара]
\label{def:mcnemar-h0}
$H_0$: вероятности дискордантных исходов равны, т.\,е.
\begin{equation}
    P(\text{$A$ верно, $B$ неверно}) = P(\text{$B$ верно, $A$ неверно}).
    \label{eq:mcnemar-h0}
\end{equation}
\end{definition}

При малых значениях $n_{10} + n_{01}$ (что типично для задач с
ограниченными тестовыми выборками) используется точный биномиальный
тест. При справедливости $H_0$ случайная величина $n_{10}$ имеет
биномиальное распределение:
\begin{equation}
    n_{10} \sim \mathrm{Bin}\!\left(n_{10} + n_{01},\; \tfrac{1}{2}\right).
    \label{eq:mcnemar-binom}
\end{equation}
$p$-значение вычисляется как вероятность наблюдения столь же или более
экстремального результата:
\begin{equation}
    p = 2\,\min\!\Bigl(
        P\bigl(X \leqslant n_{10}\bigr),\;
        P\bigl(X \geqslant n_{10}\bigr)
    \Bigr),
    \quad X \sim \mathrm{Bin}\!\left(n_{10}+n_{01},\; \tfrac{1}{2}\right).
    \label{eq:mcnemar-pvalue}
\end{equation}

При больших $n_{10} + n_{01}$ может быть использовано асимптотическое
приближение со статистикой
\begin{equation}
    \chi^{2}_{\mathrm{McN}} = \frac{(n_{10} - n_{01})^{2}}{n_{10} + n_{01}},
    \label{eq:mcnemar-chi2}
\end{equation}
имеющей при $H_0$ распределение $\chi^{2}(1)$.

%% -------------------------------------------------------------------
\subsubsection{Поправка Бонферрони}
\label{subsubsec:bonferroni}

При проведении $m$ попарных сравнений возникает проблема множественных
проверок: вероятность хотя бы одного ложноположительного результата
(family-wise error rate, FWER) возрастает с ростом $m$.

\begin{definition}[FWER]
\label{def:fwer}
Групповая вероятность ошибки первого рода определяется как
\begin{equation}
    \mathrm{FWER} = P\bigl(\text{хотя бы одно ложное отвержение } H_0\bigr).
    \label{eq:fwer}
\end{equation}
\end{definition}

Поправка Бонферрони~\cite{bonferroni1936teoria, dunn1961multiple} является
простейшим способом контроля FWER. При проведении $m$ одновременных проверок
скорректированный уровень значимости каждого отдельного теста полагается
равным
\begin{equation}
    \alpha_{\mathrm{adj}} = \frac{\alpha}{m},
    \label{eq:bonferroni}
\end{equation}
где $\alpha$ --- желаемый уровень FWER. Неравенство Бонферрони
гарантирует, что
\begin{equation}
    \mathrm{FWER} = P\!\left(\bigcup_{i=1}^{m} A_i\right)
    \leqslant \sum_{i=1}^{m} P(A_i)
    = m \cdot \alpha_{\mathrm{adj}} = \alpha,
    \label{eq:bonferroni-proof}
\end{equation}
где $A_i$ --- событие ложного отвержения $H_0$ в $i$-м тесте.

\emph{Правило применения.} Каждое из $m$ попарных сравнений
проводится критерием Макнемара, после чего его $p$-значение сравнивается
не с исходным $\alpha$, а с $\alpha_{\mathrm{adj}} = \alpha / m$;
эквивалентно, нулевая гипотеза для отдельного теста отвергается, если
$p \leqslant \alpha_{\mathrm{adj}}$, либо, что то же самое, если
скорректированное значение $p_{\mathrm{adj}} = \min(1,\; m\cdot p)$
не превосходит~$\alpha$.

В нашем случае при сравнении $k$ классификаторов попарно число
сравнений составляет $m = \binom{k}{2} = \frac{k(k-1)}{2}$. Для $k = 12$
получаем $m = 66$ и $\alpha_{\mathrm{adj}} = 0{,}05 / 66 \approx 0{,}000758$.

Поправка Бонферрони консервативна, особенно при положительной корреляции
между тестами (что характерно для попарных сравнений на одной и той же
выборке): фактический FWER, как правило, оказывается строго меньше~$\alpha$,
а статистическая мощность снижается. Менее консервативные альтернативы~---
например, последовательная процедура Холма~\cite{holm1979simple}~--- также
контролируют FWER на уровне~$\alpha$, но обладают большей мощностью.
В настоящей работе мы сознательно выбираем процедуру Бонферрони, поскольку
основной вывод формулируется в форме <<модель~A значимо превосходит
модель~B>>, и для такого утверждения предпочтительнее минимизировать риск
ложного отвержения, нежели терять часть истинных различий.


%% ===================================================================
\subsection{Бутстрап доверительные интервалы}
\label{subsec:bootstrap}

Для оценки неопределённости метрик качества классификации применяется
метод непараметрического бутстрапа~\cite{efron1979bootstrap}, не
требующий предположений о параметрическом виде распределения метрики.

Пусть $\mathcal{D} = \{(x_i, y_i, \hat{y}_i)\}_{i=1}^{n}$ ---
тестовая выборка с истинными и предсказанными метками. Пусть
$\theta = T(\mathcal{D})$ --- значение интересующей метрики
(например, $F_1$, error rate), вычисленное на выборке $\mathcal{D}$.

\begin{definition}[Бутстрап-выборка]
\label{def:bootstrap-sample}
Бутстрап-выборкой $\mathcal{D}^{*}_b$ ($b = 1, \ldots, B$) называется
выборка объёма $n$, полученная из $\mathcal{D}$ случайным выбором с
возвращением:
\begin{equation}
    \mathcal{D}^{*}_b = \bigl\{(x_{i_1}, y_{i_1}, \hat{y}_{i_1}),\;
    \ldots,\; (x_{i_n}, y_{i_n}, \hat{y}_{i_n})\bigr\},
    \quad i_1, \ldots, i_n \stackrel{\mathrm{iid}}{\sim}
    \mathrm{Uniform}\{1, \ldots, n\}.
    \label{eq:bootstrap-sample}
\end{equation}
\end{definition}

Алгоритм построения бутстрап доверительного интервала перцентильным
методом состоит из следующих шагов.

\begin{enumerate}
    \item Для каждого $b = 1, \ldots, B$ сформировать бутстрап-выборку
          $\mathcal{D}^{*}_b$ и вычислить
          $\theta^{*}_b = T(\mathcal{D}^{*}_b)$.
    \item Упорядочить бутстрап-оценки:
          $\theta^{*}_{(1)} \leqslant \theta^{*}_{(2)} \leqslant \ldots
          \leqslant \theta^{*}_{(B)}$.
    \item Построить доверительный интервал уровня $(1-\alpha)$ как
\end{enumerate}
\begin{equation}
    \mathrm{CI}_{1-\alpha} = \Bigl[\theta^{*}_{(\lfloor B\alpha/2 \rfloor)},\;
    \theta^{*}_{(\lceil B(1-\alpha/2) \rceil)}\Bigr].
    \label{eq:bootstrap-ci}
\end{equation}

В настоящей работе используется $B = 1000$ бутстрап-выборок и уровень
доверия $1 - \alpha = 0{,}95$, что обеспечивает построение $95\%$
доверительных интервалов для всех исследуемых метрик.

Обоснование состоятельности перцентильного метода основано на принципе
подстановки (plug-in principle): эмпирическое распределение
$\hat{F}_n$ аппроксимирует истинное распределение $F$, а распределение
$T(\mathcal{D}^{*}) - T(\mathcal{D})$ аппроксимирует распределение
$T(\mathcal{D}) - \theta$, где $\theta$ --- истинное значение метрики.
Состоятельность аппроксимации $\hat{F}_n \to F$ при $n \to \infty$
обеспечивает классическая теорема Гливенко--Кантелли~\cite{glivenko1933sulla, cantelli1933sulla}:
\[
    \sup_{x \in \mathbb{R}} \bigl|\hat{F}_n(x) - F(x)\bigr|
    \xrightarrow{\text{п.\,н.}} 0,
\]
а состоятельность бутстрап-оценки распределения статистики $T$ при
выполнении условий регулярности (липшицевость, конечность дисперсии)
показана в~\cite{bickel1981asymptotic, efron1993bootstrap}. Тестовые
выборки настоящей работы ($n = 500$ и $n = 250$) удовлетворяют этим
условиям для рассматриваемых метрик.
